Existence and (anti)multiplicativity: the grade involution is $\Cl(-1)$ by \claimref{std-algebra.functoriality}; reversion is the homomorphism into the opposite algebra given by \claimref{std-algebra.existence-uniqueness}; conjugation is their composite. Each is an involution because its square is a homomorphism (or, for the anti-automorphisms, the square of an anti-homomorphism is a homomorphism) fixing $V$, hence the identity by uniqueness in the universal property. They commute for the same reason: both composites are homomorphisms or anti-homomorphisms agreeing on $V$.

Signs: by \claimref{std-algebra.grade-decomposition}, $\Cl^{k}$ is spanned by products $v_1 \cdots v_k$ of mutually orthogonal vectors. The grade involution gives $(-1)^{k}$. Reversion gives $v_k \cdots v_1$, and reordering anticommuting factors takes $k(k-1)/2$ transpositions. Conjugation multiplies the two signs: $(-1)^{k + k(k-1)/2} = (-1)^{k(k+1)/2}$.
